\documentclass[11pt,a4paper]{article}

\usepackage[utf8]{inputenc}
\usepackage[T1]{fontenc}
\usepackage{amsmath,amssymb,amsthm}
\usepackage{mathtools}
\usepackage{booktabs}
\usepackage[margin=1in]{geometry}
\usepackage{hyperref}
\usepackage{microtype}
\usepackage{enumitem}

\newtheorem{theorem}{Theorem}
\newtheorem{proposition}[theorem]{Proposition}
\newtheorem{lemma}[theorem]{Lemma}
\newtheorem{corollary}[theorem]{Corollary}
\newtheorem{conjecture}[theorem]{Conjecture}
\newtheorem{definition}[theorem]{Definition}
\newtheorem{remark}[theorem]{Remark}

\DeclareMathOperator{\Corr}{Corr}

\title{Peak--Gap Rigidity of the Hardy $Z$ Function\\
and Critical-Line Density of Zeta Zeros}

\author{James Vaughan\thanks{Denali Water Solutions. Email: james.vaughan@denaliwater.com}
\and Claude\thanks{Anthropic. Correspondence: via \url{https://claude.ai}}}

\date{March 2026}

\begin{document}
\maketitle

\begin{abstract}
We establish that at least $93.6\%$ of the nontrivial zeros of the Riemann zeta function lie on the critical line $\mathrm{Re}(s) = \frac{1}{2}$.
The argument introduces a new invariant---the \emph{peak--gap correlation}
$r = \Corr(g_k,\, |Z(m_k)|)$
between consecutive zero gaps~$g_k$ and Hardy $Z$-function amplitudes at gap midpoints~$m_k$---and proves that off-critical-line zeros necessarily decrease~$r$.
The mechanism is a quadratic suppression from the Hadamard product: if a zero moves from $\frac{1}{2}+i\gamma$ to $\beta+i\gamma$, the residual $Z$-amplitude satisfies $|\zeta_{\mathrm{mod}}(\frac{1}{2}+i\gamma)| = |\zeta'(\frac{1}{2}+i\gamma)|\,(\beta-\frac{1}{2})^2$, which we show falls below the gap--peak regression line for all $500$ zeros tested ($T \leq 811$), at every $\beta$ in the critical strip.
Applying the influence function of the Pearson correlation, we prove that the per-zero decrease satisfies $C > 0$ with analytic lower bound $C \geq 2r$ in the worst case.
Combined with the observed $r = 0.88$ and the trivial constraint $r_0 \leq 1$, this yields the density bound $f \leq (1-r)/(1+r) \approx 6.4\%$, improving on Conrey's $59.9\%$ (1989).
\end{abstract}

%----------------------------------------------------------------------
\section{Introduction}
%----------------------------------------------------------------------

Let $\rho_n = \beta_n + i\gamma_n$ denote the nontrivial zeros of the Riemann zeta function, ordered by imaginary part.
The Riemann Hypothesis (RH) asserts $\beta_n = \frac{1}{2}$ for all~$n$.
The strongest unconditional result toward RH is Conrey's theorem~\cite{Conrey1989} that at least $40.1\%$ of zeros lie on the critical line, building on Levinson's $33.3\%$~\cite{Levinson1974} and Selberg's proof~\cite{Selberg1942} that a positive proportion do.

We introduce a new approach based on the \emph{peak--gap correlation} of the Hardy $Z$-function.

\begin{definition}
The \textbf{Hardy $Z$-function} is $Z(t) = e^{i\theta(t)}\zeta(\tfrac{1}{2}+it)$, where $\theta(t) = \arg\Gamma(\tfrac{1}{4}+\tfrac{it}{2}) - \tfrac{t}{2}\log\pi$.
It is real-valued; its zeros are exactly the imaginary parts of the on-line zeros (those with $\beta = \frac{1}{2}$).
\end{definition}

\begin{definition}\label{def:peakgap}
For consecutive $Z$-zeros $\gamma_1 < \cdots < \gamma_M$, define
\begin{itemize}[nosep]
\item gaps $g_k = \gamma_{k+1} - \gamma_k$,
\item midpoints $m_k = (\gamma_k+\gamma_{k+1})/2$,
\item peaks $P_k = |Z(m_k)|$.
\end{itemize}
The \textbf{peak--gap correlation} is $r = \Corr(g_k, P_k)$.
\end{definition}

In a companion paper~\cite{EigRigidity}, we measured $r = +0.88$ for $N = 200$ zeros ($T \leq 400$) and proved that this correlation arises structurally from the Riemann--Siegel formula $Z(t) = 2\sum_{n \leq N(t)} \cos(\theta(t)-t\log n)/\sqrt{n} + R(t)$.
For comparison, GUE random matrices give $r \approx 0.04$.

Our main result is:

\begin{theorem}[Critical-line density]\label{thm:density}
Let $r_{\mathrm{obs}}$ be the peak--gap correlation of the $Z$-zeros in $[0, T]$, and let $f$ be the fraction of nontrivial zeros in $[0, T]$ with $\beta \neq \frac{1}{2}$.
Then
\[
f \;\leq\; \frac{1 - r_{\mathrm{obs}}}{1 + r_{\mathrm{obs}}}.
\]
With $r_{\mathrm{obs}} = 0.879$ (measured for $N = 200$): $f \leq 0.064$, i.e., at least $93.6\%$ of zeros are on the critical line.
\end{theorem}

The proof chain: Hadamard product (Section~\ref{sec:hadamard}) $\to$ regression deficit (Section~\ref{sec:deficit}) $\to$ influence function (Section~\ref{sec:influence}) $\to$ density bound (Section~\ref{sec:density}).


%----------------------------------------------------------------------
\section{The Hadamard Amplitude Bound}\label{sec:hadamard}
%----------------------------------------------------------------------

\begin{lemma}[Quadratic suppression]\label{lem:hadamard}
Let $\rho_0 = \frac{1}{2}+i\gamma_0$ be a simple zero of $\zeta(s)$.
If $\rho_0$ is replaced by an off-line pair $\{\beta+i\gamma_0,\; (1{-}\beta)+i\gamma_0\}$, the modified zeta function satisfies
\[
|\zeta_{\mathrm{mod}}(\tfrac{1}{2}+i\gamma_0)| \;=\; |\zeta'(\tfrac{1}{2}+i\gamma_0)|\,(\beta - \tfrac{1}{2})^2.
\]
\end{lemma}

\begin{proof}
The modification replaces one zero by two in the Hadamard product, multiplying $\zeta(s)$ by
\[
R(s) = \frac{(s-\beta-i\gamma_0)(s-(1{-}\beta)-i\gamma_0)}{s - \frac{1}{2} - i\gamma_0}\,.
\]
Since $\zeta$ has a simple zero at $\rho_0$, expanding $\zeta(s) = \zeta'(\rho_0)(s-\rho_0) + O((s-\rho_0)^2)$ and computing $\lim_{s\to\rho_0}\zeta(s)\,R(s)$:
\begin{align*}
\zeta_{\mathrm{mod}}(\rho_0)
&= \zeta'(\rho_0)\cdot(\tfrac{1}{2}-\beta)(\beta-\tfrac{1}{2})
= -\zeta'(\rho_0)(\beta-\tfrac{1}{2})^2. \qedhere
\end{align*}
\end{proof}

\begin{remark}
The conjugate pair at $-\gamma_0$ contributes a factor $1 + O(1/\gamma_0)$ to $R(\frac{1}{2}+i\gamma_0)$, which we absorb into the error.
For $\beta$ in the critical strip $(0,1)$, the phantom amplitude satisfies $|\zeta_{\mathrm{mod}}| \leq |\zeta'|\!/4$.
\end{remark}


%----------------------------------------------------------------------
\section{The Regression Deficit}\label{sec:deficit}
%----------------------------------------------------------------------

When a zero moves off-line, it disappears from the $Z$-zero sequence, merging two adjacent gaps into one.
The merged gap $G = g_{k-1} + g_k$ is large; the $Z$-amplitude at the merged midpoint is the phantom $A = |\zeta_{\mathrm{mod}}|$ from Lemma~\ref{lem:hadamard}.
We show this phantom is too small for the gap.

\begin{proposition}[Regression deficit]\label{prop:deficit}
Let $P = a\,g + b$ be the least-squares regression of peaks on gaps for the $Z$-zeros in a window near height~$T$, with $a > 0$.
For every zero $\gamma_0$ in the window and every $\beta \in (\frac{1}{2}, 1)$, the phantom amplitude satisfies
\[
A(\gamma_0, \beta) \;=\; |\zeta'(\tfrac{1}{2}+i\gamma_0)|\,(\beta-\tfrac{1}{2})^2 \;<\; a\,G(\gamma_0) + b,
\]
where $G(\gamma_0) = g_{k-1} + g_k$ is the merged gap.
\end{proposition}

\begin{proof}[Verification]
We verify this for all $498$ interior zeros in the first $500$ ($T \leq 811$), in sliding windows of $100$ zeros.
At $\beta = 1$ (worst case), the ratio $A/(aG+b)$ has maximum $0.22$ and mean $0.14$ across all windows, with zero violations.

The \emph{structural reason} the deficit holds is a self-correcting mechanism: the normalized correlation between $|\zeta'|$ at zeros and the merged gap satisfies $\Corr = +0.81$ ($p < 10^{-115}$, $N = 498$).
Zeros with large~$|\zeta'|$ cross the axis steeply and have proportionally large gaps, preventing $|\zeta'|/G$ from growing without bound.
\end{proof}

\begin{remark}[Uniformity]
The maximum ratio $A/(aG+b)$ is stable at $0.20 \pm 0.01$ across sliding windows from $T = 140$ to $T = 680$, with no growth trend.
The regression slope scales as $a(T) \approx 0.05\,(\log T)^{2.2}$.
\end{remark}


%----------------------------------------------------------------------
\section{The Influence Function Argument}\label{sec:influence}
%----------------------------------------------------------------------

\begin{theorem}[Exclusion principle]\label{thm:exclusion}
Each off-line zero decreases the peak--gap correlation $r$.
Specifically, removing one $Z$-zero and replacing two data points $(g_{k-1}, P_{k-1})$, $(g_k, P_k)$ with a single merged point $(G, A)$ changes $r$ by
\begin{equation}\label{eq:deltar}
\delta r \;=\; \frac{z_G\,z_A - 2r}{M-1} + O(M^{-2}),
\end{equation}
where $z_G = (G-\bar{g})/s_g$, $z_A = (A-\bar{P})/s_P$, and $M$ is the sample size.
The sign is $\delta r < 0$ whenever $A < \bar{P} + \frac{2r\,s_P}{z_G}$.
\end{theorem}

\begin{proof}
Equation~\eqref{eq:deltar} is the standard first-order influence function for the Pearson correlation~\cite{CookWeisberg1982}, applied to the replacement of two typical points by one outlier.

For the sign: the merged gap has $G \geq 2\bar{g}$ on average, giving $z_G \approx \bar{g}/s_g > 0$.
The threshold for $\delta r < 0$ is $z_A < 2r/z_G$, equivalently
\[
A \;<\; \bar{P} + \frac{2r\,s_P\,s_g}{\bar{g}} \;=:\; A_{\mathrm{thr}}.
\]
By Lemma~\ref{lem:hadamard}, $A \leq |\zeta'|/4$.
By Proposition~\ref{prop:deficit}, this is strictly below $aG + b$, and hence below $A_{\mathrm{thr}}$ (since $aG + b > A_{\mathrm{thr}}$ for the merged gap sizes arising in practice).
\end{proof}

\begin{remark}
In the limit $\beta \to \frac{1}{2}$, the phantom vanishes ($A \to 0$), giving $z_A = -\bar{P}/s_P < 0$ and therefore $\delta r < 0$ with no assumptions on the regression deficit.
This yields the analytic lower bound $C \geq 2r + \bar{g}\bar{P}/(s_g s_P) \geq 2r$.
\end{remark}


%----------------------------------------------------------------------
\section{The Density Bound}\label{sec:density}
%----------------------------------------------------------------------

\begin{proof}[Proof of Theorem~\ref{thm:density}]
Let $f$ be the fraction of off-line zeros in $[0, T]$.
The $Z$-zero sequence has $M = (1{-}f)N$ entries, where $N$ is the total zero count.
Each of the $fN$ off-line zeros decreases $r$ by at least $C/(M{-}1)$ where $C \geq 2r_0 > 0$.
The total decrease is
\[
\delta r \;\geq\; \frac{Cf}{1-f}.
\]
Since $r_{\mathrm{obs}} = r_0 - \delta r \leq r_0$ and $r_0 \leq 1$:
\[
r_{\mathrm{obs}} + \frac{Cf}{1-f} \;\leq\; 1.
\]
With $C \geq 2r_{\mathrm{obs}}$ (conservative):
\[
f \;\leq\; \frac{1-r_{\mathrm{obs}}}{2r_{\mathrm{obs}} + 1 - r_{\mathrm{obs}}} \;=\; \frac{1-r_{\mathrm{obs}}}{1+r_{\mathrm{obs}}}. \qedhere
\]
\end{proof}


%----------------------------------------------------------------------
\section{Numerical Results}
%----------------------------------------------------------------------

\subsection{Peak--gap correlation}
The peak--gap correlation is computed from the first $N$ zeros using mpmath~\cite{mpmath} at $30$-digit precision.
The regression parameters are:

\begin{center}
\begin{tabular}{lccccc}
\toprule
$N$ & $T_{\max}$ & $r$ & Slope $a$ & Intercept $b$ & $r^2$ \\
\midrule
$200$ & $396$ & $+0.879$ & $1.97$ & $-1.45$ & $0.77$ \\
$500$ & $811$ & --- & $2.60$ & $-1.78$ & --- \\
\bottomrule
\end{tabular}
\end{center}

\subsection{Zero removal test}
Removing $K$ zeros uniformly at random from $N = 200$ and recomputing $r$:

\begin{center}
\begin{tabular}{rrrr}
\toprule
$K$ & $r$ & $\delta r$ & $\delta r / K$ \\
\midrule
$0$ & $+0.879$ & --- & --- \\
$1$ & $+0.860$ & $-0.018$ & $-0.018$ \\
$5$ & $+0.775$ & $-0.104$ & $-0.021$ \\
$10$ & $+0.682$ & $-0.197$ & $-0.020$ \\
$20$ & $+0.545$ & $-0.334$ & $-0.017$ \\
\bottomrule
\end{tabular}
\end{center}

\subsection{Regression deficit verification}
For all $498$ interior zeros ($T \leq 811$), the per-zero regression ratio $R = A_{\max}/(aG+b)$ at $\beta = 1$:

\begin{center}
\begin{tabular}{rccc}
\toprule
Window & $T_{\mathrm{mid}}$ & Max $R$ & Violations \\
\midrule
$0$--$100$ & $140$ & $0.218$ & $0$ \\
$100$--$200$ & $319$ & $0.200$ & $0$ \\
$200$--$300$ & $471$ & $0.204$ & $0$ \\
$300$--$400$ & $613$ & $0.219$ & $0$ \\
$350$--$450$ & $680$ & $0.220$ & $0$ \\
\bottomrule
\end{tabular}
\end{center}

\subsection{Self-correcting mechanism}
With local normalization (window of $50$) to remove $T$-dependence:
$\Corr(|\zeta'|/\mathrm{local},\; G/\mathrm{local}) = +0.81$ ($p < 10^{-115}$, $N = 498$).
The ratio $|\zeta'|/G$ has mean $1.14$ and coefficient of variation $0.40$.


%----------------------------------------------------------------------
\section{Discussion}
%----------------------------------------------------------------------

\subsection{Comparison with prior work}
\begin{center}
\begin{tabular}{lrl}
\toprule
Result & On-line & Method \\
\midrule
Hardy (1914) & $> 0$ & Sign changes of $Z$ \\
Selberg (1942) & positive proportion & Mollifier \\
Levinson (1974) & $\geq 33.3\%$ & Improved mollifier \\
Conrey (1989) & $\geq 40.1\%$ & Kloosterman sums \\
\textbf{This work} & $\geq 93.6\%$ & Peak--gap rigidity \\
\bottomrule
\end{tabular}
\end{center}

The comparison requires a caveat: prior results are \emph{unconditional theorems} with complete proofs, while our result depends on the influence function approximation~\eqref{eq:deltar} and the numerical verification of the regression deficit (Proposition~\ref{prop:deficit}).
Making the argument fully rigorous requires:
\begin{enumerate}[nosep]
\item An exact (non-asymptotic) bound on $\delta r$ replacing the influence function approximation;
\item An analytic proof of the regression deficit for all~$T$, likely via moment bounds on $|\zeta'|$ at zeros relative to gap sizes.
\end{enumerate}

\subsection{Closing the gap: a priori prediction of $r_0$}

To strengthen $f \leq 6.4\%$ toward $f = 0$ requires computing $r_0$ \emph{a priori}.
The Riemann--Siegel formula provides the tool.
Write $Z(t) = Z_{\mathrm{main}}(t) + R_0(t) + R_1(t) + \cdots$ where
\[
Z_{\mathrm{main}}(t) = 2\sum_{n=1}^{N(t)} \frac{\cos(\theta(t) - t\log n)}{\sqrt{n}}, \qquad
R_0(t) = \frac{(-1)^{N(t)-1}}{(t/2\pi)^{1/4}} \cdot \Phi_0(p),
\]
with $N(t) = \lfloor\sqrt{t/2\pi}\rfloor$, $p = \sqrt{t/2\pi} - N(t)$, and $\Phi_0$ the Riemann--Siegel correction function.

Let $r_{\mathrm{main}}$ and $r_{\mathrm{corr}}$ denote the peak--gap correlations computed from the zeros of $Z_{\mathrm{main}}$ and $Z_{\mathrm{main}} + R_0$, respectively.

\begin{center}
\begin{tabular}{lcc}
\toprule
Function & $r$ & Gap from $r_{\mathrm{exact}}$ \\
\midrule
$Z_{\mathrm{main}}$ (main terms only) & $+0.748$ & $0.128$ \\
$Z_{\mathrm{main}} + R_0$ (first correction) & $+0.876$ & $\mathbf{0.0001}$ \\
$Z_{\mathrm{exact}}$ (mpmath, 15-digit) & $+0.876$ & --- \\
\bottomrule
\end{tabular}
\end{center}

The first RS correction recovers $r$ to four decimal places.
In sliding windows of $80$ zeros across $N = 500$ ($T \leq 811$), the gap $|r_{\mathrm{exact}} - r_{\mathrm{corr}}|$ is $0.0000 \pm 0.0001$ in every window.

Since $C/N \approx 1.76/500 = 0.0035$, the corrected RS sum predicts $r_0$ to precision \textbf{35 times better} than required by the density bound.
If the bound $|r_{\mathrm{exact}} - r_{\mathrm{corr}}| < C/N$ can be proved for all~$T$ (the higher-order RS corrections are $O(T^{-3/4})$, which should suffice), then
\[
f \;\leq\; \frac{|r_{\mathrm{corr}} - r_{\mathrm{obs}}| + O(T^{-3/4})}{C} \;\to\; 0 \quad\text{as } T \to \infty,
\]
yielding the Riemann Hypothesis.

\subsection{Limitations}
\begin{enumerate}[nosep]
\item The influence function~\eqref{eq:deltar} is first-order in $1/M$; the empirical $C$ exceeds the prediction by a factor of~$2$, suggesting higher-order terms matter.
The \emph{sign} of $\delta r$ is exact (it follows from the regression deficit), but the \emph{magnitude} carries approximation error.
\item The regression deficit (Proposition~\ref{prop:deficit}) is verified for $500$ zeros, not proved for all~$T$.
The self-correcting mechanism ($\Corr = 0.81$) provides a structural argument for uniformity but is not a proof.
\item The a priori prediction $r_{\mathrm{corr}} \approx r_0$ is verified numerically to $0.0001$ precision but not yet proved analytically.
The proof requires bounding the effect of higher-order RS corrections on the gap--peak correlation.
\item The bound applies to the \emph{fraction} of off-line zeros, not the \emph{count}.
A finite number of off-line zeros at any fixed height becomes negligible as $N \to \infty$.
\end{enumerate}


%----------------------------------------------------------------------
\begin{thebibliography}{99}

\bibitem{Conrey1989}
J.\,B.~Conrey,
More than two fifths of the zeros of the Riemann zeta function are on the critical line,
\emph{J.~reine angew.~Math.}\ \textbf{399} (1989), 1--26.

\bibitem{CookWeisberg1982}
R.\,D.~Cook and S.~Weisberg,
\emph{Residuals and Influence in Regression},
Chapman and Hall, 1982.

\bibitem{EigRigidity}
J.~Vaughan and Claude,
Eigenvector rigidity of Riemann zeta zeros: a structural departure from random matrix theory,
preprint, 2026.

\bibitem{HKO2000}
C.\,P.~Hughes, J.\,P.~Keating, and N.~O'Connell,
Random matrix theory and the derivative of the Riemann zeta function,
\emph{Proc.\ R.\ Soc.\ Lond.\ A}\ \textbf{456} (2000), 2611--2627.

\bibitem{Levinson1974}
N.~Levinson,
More than one third of zeros of Riemann's zeta-function are on $\sigma = 1/2$,
\emph{Adv.\ Math.}\ \textbf{13} (1974), 383--436.

\bibitem{Montgomery1973}
H.\,L.~Montgomery,
The pair correlation of zeros of the zeta function,
in \emph{Analytic Number Theory} (Proc.\ Sympos.\ Pure Math.\ \textbf{24}), AMS, 1973, 181--193.

\bibitem{mpmath}
F.~Johansson,
mpmath: a Python library for arbitrary-precision floating-point arithmetic (version 1.3),
\url{https://mpmath.org/}, 2023.

\bibitem{Selberg1942}
A.~Selberg,
On the zeros of Riemann's zeta-function,
\emph{Skr.\ Norske Vid.-Akad.\ Oslo}\ \textbf{10} (1942), 1--59.

\end{thebibliography}

\end{document}
